%% kapa: exam=1 round=23 year=2012 period=2 form=A subject=appraisal_law subject_ko=부동산관계법규 seq=25 number=65
%% source: Q-Net 감정평가사 기출문제 (cst00309) artlSeq=5052580 file=raw/1008/5052580/2012년 제23회 감정평가사 제1차시험 기출문제(2교시 A형).pdf
%% conversion: kordoc 4.24.51 -> tools/gen1.py (auto); review=proofread-sonnet; flags=none
\begin{question}{65}
국유재산법령상 국유재산에 관한 설명으로 옳지 않은 것은?
\begin{choices}
  \item 중앙관서의 장이 국유재산의 관리에 관련된 법령을 개정하려면 그 내용에 관하여 총괄청 및 감사원과 협의하여야 한다.
  \item 총괄청이나 중앙관서의 장은 소유자 없는 부동산을 국유재산으로 취득한다.
  \item 확정판결에 의한 경우에는 일반재산에 대하여 사권을 설정할 수 있다.
  \item 위법한 재산관리로 사용승인이 철회되어 총괄청에 인계된 행정재산도 그 용도가 폐지되는 것은 아니다.
  \item 총괄청은 일반재산을 보존용재산으로 전환하여 관리할 수 있다.
\end{choices}
\end{question}
\answer{4}
